\documentclass[handout]{beamer}
\mode<presentation>
\usepackage{xcolor}
\usepackage[toc]{multitoc}
\usepackage{color}
\usepackage{graphicx}
\graphicspath{{../buku_ajar/buku_ajar_logika_matematika/figures/}}
\usepackage{tikz}
\usetikzlibrary{shapes,arrows,positioning,calc}
\usepackage{listings}
\usepackage{lmodern}
\usetheme[secheader]{Boadilla}
\usecolortheme{whale}
\setbeamertemplate{caption}[numbered]
\setbeamertemplate{bibliography item}[text]
\setbeamertemplate{navigation symbols}{}

\theoremstyle{definition}
\newtheorem{definisi}[theorem]{Definisi}
\theoremstyle{example}
\newtheorem{contoh}[theorem]{Contoh}

\renewcommand*{\multicolumntoc}{2}

\setbeamertemplate{footline}
{
  \leavevmode%
  \hbox{%
  \begin{beamercolorbox}[wd=.333333\paperwidth,ht=2.25ex,dp=1ex,center]{author in head/foot}%
    \usebeamerfont{author in head/foot}\insertshortauthor%
  \end{beamercolorbox}%
  \begin{beamercolorbox}[wd=.433333\paperwidth,ht=2.25ex,dp=1ex,center]{title in head/foot}%
    \usebeamerfont{title in head/foot}\insertshorttitle
  \end{beamercolorbox}%
  \begin{beamercolorbox}[wd=.233333\paperwidth,ht=2.25ex,dp=1ex,right]{date in head/foot}%
    \usebeamerfont{date in head/foot}\insertshortdate{}\hspace*{2em} \insertframenumber{} / \inserttotalframenumber\hspace*{2ex} 
  \end{beamercolorbox}}%
  \vskip0pt%
}

\subtitle[Logika Matematika]{Logika Matematika}
\title[Logika Predikat]{Logika Predikat, Fungsi Proposisional, dan Universe of Discourse}
\author{Cecep Suwanda, S.Si., M.Kom.}
\date{Maret 2025}
\institute{Universitas Bale Bandung}

\begin{document}

\maketitle

\section*{Daftar Isi}
\begin{frame}
\frametitle{Daftar Isi}
\tableofcontents
\end{frame}

% ============================================================
% SECTION 6-1: Keterbatasan Logika Proposisi
% ============================================================
\section{Keterbatasan Logika Proposisi}

\begin{frame}
\frametitle{Argumen Silogisme Klasik}
\begin{block}{Argumen Aristoteles}
\textbf{Premis 1:} Semua manusia pasti akan mati.\\
\textbf{Premis 2:} Socrates adalah manusia.\\
\textbf{Kesimpulan:} Socrates pasti akan mati.
\end{block}
Secara intuitif valid. Uji dengan logika proposisi?
\end{frame}

\begin{frame}
\frametitle{Mengapa Logika Proposisi Gagal?}
Representasi atomik: $p$ = semua manusia mati, $q$ = Socrates manusia, $r$ = Socrates mati.\\
Struktur: $p, q \vdash r$
\begin{block}{Keterbatasan}
Logika proposisi memperlakukan setiap kalimat sebagai \textit{black box}---struktur internal (subjek--predikat) tidak dianalisis. Tidak ada hubungan inferensial langsung antara $p$, $q$, $r$.
\end{block}
\end{frame}

\begin{frame}
\frametitle{Logika Predikat (FOL)}
\begin{definisi}[Logika Predikat / First-Order Logic]
Perluasan logika proposisional yang memungkinkan penguraian kalimat menjadi \textbf{subjek} (objek) dan \textbf{predikat} (sifat/relasi). Memperkenalkan variabel dan kuantifikasi.
\end{definisi}
Representasi: $\text{Manusia}(\text{Socrates})$, $\text{LebihBesar}(5, 3)$.
\end{frame}

% ============================================================
% SECTION 6-2: Term, Konstanta, Variabel
% ============================================================
\section{Struktur: Term, Konstanta, Variabel}

\begin{frame}
\frametitle{Definisi Term}
\begin{definisi}[Term]
Elemen/entitas yang menjadi subjek pernyataan, sifat, atau relasi. Berperan sebagai \textbf{kata benda}; tidak bernilai benar/salah.
\end{definisi}
Tiga kategori: \textbf{konstanta}, \textbf{variabel}, hasil \textbf{aplikasi fungsi} terhadap term.
\end{frame}

\begin{frame}
\frametitle{Konstanta dan Variabel}
\begin{block}{Konstanta}
Penamaan spesifik objek: ``Socrates'', ``Jakarta'', $1$, $\pi$.
\end{block}
\begin{block}{Variabel}
Subjek belum ditetapkan; dapat disubstitusi. Simbol: $x$, $y$, $z$. Kata ganti ``ia'' berperan seperti variabel.
\end{block}
\end{frame}

\begin{frame}
\frametitle{Fungsi Pemetaan}
Fungsi memetakan term $\to$ term baru (bukan nilai kebenaran).
\begin{contoh}
\begin{itemize}
  \item $f(x) = x^2 + 1$; untuk $x=3$ $\Rightarrow$ term $f(3)=10$.
  \item $ayah(x)$ = ``ayah kandung dari $x$''.
  \item $dekan(F)$ = individu yang menjabat dekan fakultas $F$.
\end{itemize}
\end{contoh}
\end{frame}

% ============================================================
% SECTION 6-3: Predikat dan Arity
% ============================================================
\section{Predikat dan Arity}

\begin{frame}
\frametitle{Definisi Predikat}
\begin{definisi}[Predikat]
Atribut, klasifikasi, properti, atau relasi yang dilekatkan pada satu atau lebih term. Menentukan nilai kebenaran (T/F) bersama interpretasi.
\end{definisi}
\begin{contoh}
``Socrates adalah manusia'': $M(s)$ dengan $M(x) \equiv x$ manusia, $s \equiv$ Socrates.
\end{contoh}
\end{frame}

\begin{frame}
\frametitle{Arity: Unary, Binary, $n$-ary}
\begin{block}{Unary (1-arity)}
$G(n) \equiv n$ bilangan genap. $W(x) \equiv x$ warga Afrika.
\end{block}
\begin{block}{Binary (2-arity)}
$L(x,y) \equiv x$ lebih besar daripada $y$. Urutan argumen penting!
$L(100,2)$ = T; $L(2,100)$ = F.
\end{block}
\begin{block}{$n$-ary}
$B(x,y,z)$ = kota $x$ terletak antara kota $y$ dan $z$.
\end{block}
\end{frame}

% ============================================================
% SECTION 6-4: Fungsi Proposisional
% ============================================================
\section{Fungsi Proposisional}

\begin{frame}
\frametitle{Definisi Fungsi Proposisional}
\begin{definisi}[Fungsi Proposisional / Kalimat Terbuka]
Ekspresi dari predikat yang mengikat setidaknya satu \textbf{variabel bebas}. Secara umum \textbf{belum} bernilai benar/salah hingga variabel diinterpretasikan.
\end{definisi}
\end{frame}

\begin{frame}
\frametitle{Dari Kalimat Terbuka ke Tertutup}
\begin{contoh}
$P(x) \equiv$ ``$x$ adalah presiden pertama RI.''\\
Selama $x$ tidak ditetapkan $\Rightarrow$ kalimat terbuka.
\end{contoh}
\begin{block}{Substitusi Konstanta}
$P(a)$ dengan $a$ = B.J. Habibie $\Rightarrow$ F.\\
$P(s)$ dengan $s$ = Ir. Soekarno $\Rightarrow$ T.
\end{block}
\end{frame}

% ============================================================
% SECTION 6-5: Semesta Pembicaraan
% ============================================================
\section{Semesta Pembicaraan}

\begin{frame}
\frametitle{Universe of Discourse ($\mathcal{U}$)}
\begin{definisi}[Semesta Pembicaraan]
Himpunan objek yang \textbf{diperbolehkan} dan \textbf{relevan} sebagai nilai substitusi untuk variabel. Menetapkan domain yang sah.
\end{definisi}
Tanpa $\mathcal{U}$: ambiguitas, substitusi tidak selaras (mis. ``Gunung Merapi'' ke predikat aritmetika $\Rightarrow$ \textit{type mismatch}).
\end{frame}

\begin{frame}
\frametitle{Pengaruh Semesta}
\begin{contoh}
$M$: $x^2 \geq 0$ untuk setiap $x$.
\begin{itemize}
  \item $\mathcal{U} = \mathbb{R}$ $\Rightarrow$ True.
  \item $\mathcal{U} = \mathbb{C}$, $x = 3i$ $\Rightarrow$ $x^2 = -9$, tidak terpenuhi.
\end{itemize}
\end{contoh}
\end{frame}

% ============================================================
% SECTION 6-6: Kalimat Tertutup
% ============================================================
\section{Kalimat Tertutup (Proposisi)}

\begin{frame}
\frametitle{Dua Cara Memperoleh Kalimat Tertutup}
\begin{enumerate}
  \item \textbf{Substitusi konstanta}: semua variabel bebas diganti konstanta sesuai semesta.
  \item \textbf{Kuantifikasi}: variabel diikat dengan $\forall$ atau $\exists$.
\end{enumerate}
\end{frame}

\begin{frame}
\frametitle{Contoh}
\begin{block}{Substitusi}
$L(x,y)$ = ``$x$ berlari lebih cepat daripada $y$''.\\
$L(\text{Snoopy}, \text{Kura-Kura Ninja})$ = proposisi tertutup.
\end{block}
\begin{block}{Kuantifikasi}
$P(x)$ = ``$x$ tidak hidup abadi'' ($\mathcal{U}$ = organisme).\\
$\forall x\, P(x)$, $\exists x\, P(x)$ = proposisi tertutup.
\end{block}
\end{frame}

% ============================================================
% SECTION 6-7: Studi Kasus Pemodelan
% ============================================================
\section{Studi Kasus: Pemodelan Naratif}

\begin{frame}
\frametitle{Skenario}
\textit{``Anton mahasiswa Informatika di Universitas Nusantara. Univ. Nusantara berlokasi di Jakarta. Anton menyukai Logika Matematika. Anton bukan ketua BEM; Sinta ketua BEM. Anton IPK lebih tinggi dari Sinta. Budi sahabat Anton.''}
\end{frame}

\begin{frame}
\frametitle{Ekstraksi Term dan Predikat}
\begin{block}{Konstanta}
$a$=Anton, $b$=Budi, $s$=Sinta, $u$=Univ. Nusantara, $j$=Jakarta, $m$=Logika Matematika
\end{block}
\begin{block}{Predikat}
Unary: $M(x)$, $B(x)$. Binary: $K(x,y)$, $L(x,y)$, $S(x,z)$, $P(x,y)$, $H(x,y)$
\end{block}
\end{frame}

\begin{frame}
\frametitle{Klausa Formal}
\begin{itemize}
  \item $M(a) \wedge K(a,u)$ --- Anton mahasiswa di Univ. Nusantara
  \item $L(u,j)$ --- Univ. berlokasi di Jakarta
  \item $S(a,m)$ --- Anton suka Logika Matematika
  \item $B(s) \wedge \neg B(a)$ --- Sinta ketua BEM, Anton bukan
  \item $P(a,s)$ --- IPK Anton lebih tinggi dari Sinta
  \item $H(b,a)$ --- Budi sahabat Anton
\end{itemize}
\end{frame}

% ============================================================
% RANGKUMAN
% ============================================================
\section{Rangkuman}

\begin{frame}
\frametitle{Rangkuman}
\begin{itemize}
  \item Logika proposisi tidak menganalisis struktur internal (subjek--predikat).
  \item \textbf{Term}: konstanta, variabel, aplikasi fungsi; tidak bernilai T/F.
  \item \textbf{Predikat}: sifat/relasi; arity = banyaknya argumen.
  \item \textbf{Fungsi proposisional}: predikat + variabel bebas; belum T/F.
  \item \textbf{Universe of Discourse}: domain substitusi yang sah.
  \item Kalimat tertutup: substitusi konstanta atau kuantifikasi.
\end{itemize}
\end{frame}

\begin{frame}
\frametitle{Kata Kunci}
\begin{center}
FOL $\bullet$ Term $\bullet$ Konstanta $\bullet$ Variabel $\bullet$ Predikat $\bullet$ Arity\\
Fungsi Proposisional $\bullet$ Kalimat Terbuka $\bullet$ Universe of Discourse
\end{center}
\end{frame}

\end{document}
